\section{Предварительные требования}
\headingtextsep

Для развёртывания системы «Контур» в текущей реализации требуется подготовленная рабочая среда, позволяющая запустить многокомпонентное приложение в контейнерном режиме. Поскольку проект состоит из нескольких взаимосвязанных сервисов, предварительные требования относятся не только к операционной системе, но и к средствам контейнеризации, сетевому окружению и локальной конфигурации.

К обязательным условиям относятся:

\begin{itemize}
\item компьютер под управлением Windows, Linux или macOS, поддерживающий запуск Docker и контейнерную виртуализацию;
\item установленный Docker Engine или Docker Desktop вместе с плагином Docker Compose;
\item доступ к исходному коду репозитория и возможность использовать файл конфигурации окружения на основе заготовки .env.example;
\item свободные сетевые порты 5432, 8080 и 5173, используемые соответственно для PostgreSQL, API и web-интерфейса;
\item доступ к сети для получения базовых контейнерных образов и загрузки зависимостей при первой сборке;
\item достаточный объем дискового пространства для хранения образов, контейнеров и постоянного тома базы данных.
\end{itemize}

На уровне прикладной конфигурации система предполагает наличие четырех основных контейнеров: базы данных PostgreSQL, API-сервиса, web-клиента и MCP-моста. Дополнительно должен быть подготовлен постоянный том для хранения данных СУБД. В учебной и локальной поставке этого набора достаточно для запуска системы без внешнего оркестратора.

Следует учитывать, что рассматриваемый в отчете способ развёртывания ориентирован прежде всего на локальный и self-hosted контур первой версии. Для более строгой production-среды могут потребоваться дополнительные элементы: внешний провайдер аутентификации, выделенное объектное хранилище файлов и отдельные механизмы резервного копирования и мониторинга.
